% Evidence-safe revision of the July 2026 spectral/entropy manuscript.
% This file deliberately excludes all unverified end-to-end benchmark claims.
\pdfsuppressptexinfo=-1
\pdfinfoomitdate=1
\pdftrailerid{}

\documentclass{article}

\usepackage[preprint,nonatbib]{neurips_2026}
\usepackage[numbers,sort&compress]{natbib}
\usepackage[utf8]{inputenc}
\usepackage[T1]{fontenc}
\usepackage{amsmath,amssymb,amsfonts,amsthm}
\usepackage{booktabs}
\usepackage{enumitem}
\usepackage{microtype}
\usepackage{tabularx}
\usepackage{xcolor}
\usepackage{xspace}
\usepackage{hyperref}
\usepackage{url}

\hypersetup{
  colorlinks=true,
  linkcolor=blue!55!black,
  citecolor=blue!55!black,
  urlcolor=blue!55!black
}

\newtheorem{proposition}{Proposition}
\newtheorem{lemma}{Lemma}
\newtheorem{remark}{Remark}
\newcommand{\grpo}{\textsc{GRPO}\xspace}
\newcommand{\zvf}{\textsc{ZVF}\xspace}
\newcommand{\E}{\mathbb{E}}
\newcommand{\Prob}{\mathbb{P}}

\title{Zero Advantage Is Not Missing Information:\\
A Mechanism Audit of Trajectory-Spectral Auxiliary Advantages}

\author{%
  Arvind C R \\
  PES University \\
  \texttt{arvindcr4@gmail.com}
}

\begin{document}
\maketitle

\begin{abstract}
Group-relative policy optimization assigns zero centered advantage whenever
every completion in a sampled group receives the same reward.  We argue that
this is an absence of task feedback rather than hidden information.  No
transform placed only in the forward path can restore a policy gradient once
the centered advantage vanishes; only an explicit auxiliary advantage can,
and that changes the objective.  We dissect one such prototype: a standardized
score built from pairwise distances between low-order Legendre features of
sampled trajectories, with an adaptive Givens rotation for coordinate
elimination.  Three results are exact.  The continuous Parseval identity holds
for the reconstructed truncated expansion, while the implemented equispaced
rule is first-order.  The rotation preserves norm precisely because its
subsequent projection discards nothing, removing about $10^{-17}$ of the norm
where plain truncation removes $10^{-3}$--$10^{-2}$.  Binary verifier rewards
impose a homogeneous-group floor of $2^{1-G}$.  On synthetic fixtures the
coupled objective produces a gradient exactly proportional to its coupling
weight; this is not evidence of learning.  We preregister alignment, placebo,
paired-seed, and ablation tests and report no benchmark result.
\end{abstract}

\section{Introduction}

Group Relative Policy Optimization and related critic-free objectives compare
several completions for the same prompt and standardize their rewards within
the sampled group \citep{shao2024deepseekmath}.  If all rewards agree, the
centered reward term is exactly zero.  Recent work studies this phenomenon as
gradient starvation or advantage collapse and proposes changes to sampling,
normalization, or the objective \citep{nie2026gradient,he2026advantage}.

An equal-reward group is observationally ambiguous.  All completions may be
wrong, all may be correct, the verifier may be coarse or mistaken, or the
completions may differ in ways irrelevant to task success.  A representation
distance can reveal that trajectories differ, but cannot by itself reveal
which difference is useful.  Consequently, adding such a distance to the
advantage is a new auxiliary objective.  It is not recovery of a missing
reward, and a nonzero gradient is not evidence that the gradient is beneficial.

This paper evaluates a narrow prototype under that interpretation.  It makes
five evidence-bounded contributions:

\begin{enumerate}[leftmargin=*,itemsep=2pt]
  \item We separate a missing reward signal from an auxiliary objective: no
  forward-path transform can create an update when its centered advantage is
  zero, whereas explicit auxiliary coupling changes the optimized objective.
  \item We define a trajectory-spectral auxiliary score and locate the
  exactness boundary between continuous Legendre identities and the implemented
  equispaced discretization.
  \item We audit adaptive Givens elimination and show that its projection is
  redundant by construction: the rotation already annihilates the coordinate
  it projects away.
  \item We give the reusable binary-reward feasibility check
  $\Prob(\text{homogeneous})\geq 2^{1-G}$ and state what procedural changes are
  required for any lower reported rate.
  \item We release traceable synthetic and audit measurements and preregister
  alignment, placebo, paired-seed, equal-compute, and ablation gates.  We make
  no policy-learning claim.
\end{enumerate}

\paragraph{Evidence status, stated before any result.}
This paper contains proofs, executable modules, 16 focused unit/harness tests,
a deterministic CPU ledger, and an independent audit covering 77 recorded
shape, mask, quadrature, isometry, gradient, and feasibility cases.  It contains
no language-model fine-tuning, held-out benchmark accuracy, multi-seed
comparison, placebo condition, or calibration of the trajectory score against
an external target.  We therefore make no benchmark-accuracy,
sample-efficiency, deployment, mitigation, or superiority claim.

\section{Related Work and Scope}

\grpo removes a learned critic by constructing within-group advantages
\citep{shao2024deepseekmath}.  DAPO combines dynamic sampling with a
token-level loss design \citep{yu2025dapo}; recent analyses isolate
group-variance pathologies and propose direct remedies
\citep{nie2026gradient,he2026advantage}.  These methods motivate the problem
but do not validate the trajectory-derived auxiliary score studied here.

Orthogonal polynomial expansions are classical tools for representing
functions on compact domains.  HiPPO and Legendre Memory Units use scaled
Legendre projections for sequence memory, and later work studies the discrete
basis and parallel training choices directly
\citep{gu2020hippo,voelker2019lmu,chilkuri2021parallelizing}.  PoPE applies
Legendre polynomials to the token axis of a language model
\citep{aggarwal2024pope}.  The closest work by purpose uses within-group
hidden-state distances for span-level credit assignment and, critically,
specifies how that signal enters the advantage \citep{chen2026shear}.
Entropy-guided advantage shaping is also an existing response to zero-variance
prompts \citep{le2025noprompt}.  Our use is therefore a narrower recombination:
a masked trajectory projection supplies a group-relative auxiliary score.  We
do not claim that low modes are semantic or high modes are noise without an
external validation target.

Norm preservation through composed Givens rotations is established
\citep{dorobantu2016dizzyrnn,jing2017eunn,frerix2019givens}.  Planar rotations
already occur in transformer attention through rotary position embedding
\citep{su2024roformer}, and have been used for parameter-efficient adaptation,
compression, and transformer blocks
\citep{ma2024qoft,su2025rotatekv,abuayyash2026threephase}.  Attention-entropy
collapse is likewise a known transformer failure mode \citep{zhai2023entropy}.
The implementation studied here uses a data-dependent rotation as a
coordinate-elimination primitive, not a quantum algorithm; we therefore remove
the earlier ``quantum-inspired'' label.  We claim no novelty for its
constituents.  The contribution is a negative mechanism audit plus a proposed
trajectory-spectral coupling whose usefulness remains to be tested.

\section{Problem Statement and Causal Boundary}

For prompt $x$, let $y_1,\ldots,y_G$ be sampled completions with scalar rewards
$r_1,\ldots,r_G$.  On active rows, define
\begin{equation}
  A_i^{\mathrm{reward}}
  = \frac{r_i-\bar r}{s_r+\epsilon},
  \qquad
  \bar r=\frac{1}{G}\sum_{j=1}^G r_j.
\end{equation}

\begin{proposition}[Homogeneous rewards remove the centered reward signal]
If $r_1=\cdots=r_G$, then $A_i^{\mathrm{reward}}=0$ for every row.  Any
score-function term weighted only by these advantages is therefore zero.
\end{proposition}

\begin{proof}
Homogeneity gives $r_i=\bar r$ for every $i$, so every numerator is zero.
\end{proof}

For binary rewards with per-completion success probability $p$, the probability
that an independent group is homogeneous is
\begin{equation}
  \Prob(\text{homogeneous}\mid p,G)=p^G+(1-p)^G.
\end{equation}
This identity diagnoses when reward contrast is unavailable.  It neither
identifies the latent cause nor supplies a replacement preference.
For binary rewards, the smallest possible homogeneous-group probability over
$p$ is $2^{1-G}$, attained at $p=1/2$.  In particular, $G=4$ has a $12.5\%$
floor.  Any lower reported rate must therefore use a larger group, non-binary
rewards, dynamic rejection/resampling, or a different denominator; it cannot
be compared to this identity without stating that procedural change.

Our prototype defines a trajectory score $s_i$ and uses
\begin{equation}
  A_i^{\mathrm{aux}}
  = \frac{s_i-\bar s}{s_s+\epsilon_s},
  \qquad
  \widetilde A_i
  = A_i^{\mathrm{reward}}+\lambda A_i^{\mathrm{aux}}.
  \label{eq:combined-advantage}
\end{equation}
Equation~\eqref{eq:combined-advantage} is an objective modification.  When
rewards are homogeneous, it can create a nonzero update if the trajectory
scores differ, but its direction is justified only if $s_i$ is shown to align
with future task reward or a valid process-level target.

\section{Trajectory-Spectral Auxiliary Score}

Let $X_i\in\mathbb{R}^{L\times D}$ denote a completion feature sequence.  The
artifact uses either supplied hidden states or a two-channel fixture consisting
of token log probabilities and normalized position.  For
$t_\ell=2\ell/(L-1)-1$ and Legendre polynomials $P_n$, it computes
\begin{equation}
  \widehat c_{i,n}
  = \frac{2n+1}{L_{\mathrm{eff}}}
    \sum_{\ell=0}^{L-1} M_{i,\ell}X_{i,\ell}P_n(t_\ell),
  \label{eq:discrete-coeff}
\end{equation}
where $M$ is a completion mask.  Pairwise distance is
\begin{equation}
  d_N(i,j)=\sum_{n=0}^{N-1}
  \frac{2}{2n+1}\left\|\widehat c_{i,n}-\widehat c_{j,n}\right\|_2^2,
  \qquad
  s_i=\frac{1}{G-1}\sum_{j\ne i}d_N(i,j).
\end{equation}
Throughout, $d_N$ is a squared weighted modal distance; we never take a square
root.  This matters because rooting is a nonlinear reparameterization that
changes the relative weight of near and far pairs, the within-group variance of
$s_i$, and hence the update induced by any fixed $\lambda$.  In the artifact,
$s_i$ is detached from the policy graph before it is standardized and coupled;
a non-detached score is a different algorithm.

\subsection{What is exact, and what is approximate}

For arbitrary coefficient vectors $c_n$ define the reconstructed function
$\widehat X_N(t)=\sum_{n=0}^{N-1}c_nP_n(t)$.  Orthogonality gives the exact
identity
\begin{equation}
  \int_{-1}^{1}\left\|\widehat X_N(t)\right\|_2^2\,dt
  =\sum_{n=0}^{N-1}\frac{2}{2n+1}\|c_n\|_2^2.
  \label{eq:parseval-truncated}
\end{equation}
Equation~\eqref{eq:parseval-truncated} is exact for the reconstructed expansion.
It does not imply that Eq.~\eqref{eq:discrete-coeff} equals the continuous
projection of the unknown trajectory function.  The implemented equally
spaced sum is neither Gauss--Legendre quadrature nor an exact discrete
orthogonal transform.  Its approximation error depends on $L$, $N$, the mask,
and regularity of the interpolated sequence.

To expose this boundary, the independent audit measured three different
discretization errors at $N=8$, including the coefficient error that the
trajectory distance actually consumes.

\begin{table}[t]
\centering
\caption{Discretization errors for the implemented equispaced rule at $N=8$.
``Gram'' is the maximum off-diagonal deviation from the continuous Gram matrix;
``operator'' is the relative Frobenius deviation of the code's
analysis--synthesis operator from identity; ``coefficient'' is relative error
against Gauss--Legendre projection.  Values are deterministic float64 CPU
measurements, not learning results.}
\label{tab:gram-error}
\scriptsize
\begin{tabular}{rcccc}
\toprule
Length $L$ & Gram & Operator & Coefficient & Trapezoid coefficient \\
\midrule
8    & .530 & 4.18   & 2.77   & 1.47 \\
32   & .0901 & .745  & .462   & .0875 \\
128  & .0174 & .152  & .101   & .00526 \\
512  & .00402& .0359 & .0242  & .000325 \\
2048 & .000983& .00885 & .00600 & .0000203 \\
\bottomrule
\end{tabular}
\end{table}

Reporting only the Gram column would flatter the implementation.  At length
$2048$, the code's operator still differs from identity by $0.885\%$, while the
coefficient error is $0.600\%$.  Replacing uniform endpoint weights by
trapezoid weights reduces the latter to $0.00203\%$, a $296\times$ improvement
with unchanged asymptotic cost.  The implemented rectangle rule is first-order;
the trapezoid rule is second-order.  The basis is also inadmissible when
$N>L_{\mathrm{eff}}$: for $L=4,N=8$, the Gram matrix has rank $4$ and condition
number $1.1\times10^{19}$.  In a separate padding-grid stress test, holding the
effective tokens fixed while changing padded length produced coefficient
discrepancies as large as $153\%$.  A training implementation must construct
coordinates from active tokens, use a mask-aware discrete basis, and raise on
rank-deficient configurations.

\section{Adaptive Givens Coordinate Elimination}

For coordinates $i\ne j$, let $G(i,j,\theta)$ be a planar rotation.  With
$\theta=\operatorname{atan2}(x_j,x_i)$ and the sign convention used by the
artifact,
\begin{equation}
  (Gx)_i=\sqrt{x_i^2+x_j^2},
  \qquad
  (Gx)_j=0.
\end{equation}

\begin{lemma}[Norm preservation of the implemented elimination]
A Givens rotation preserves $\ell_2$ norm.  If its angle is selected as above,
projecting coordinate $j$ to zero after the rotation leaves the norm unchanged
in exact arithmetic because that coordinate is already zero.
\end{lemma}

\begin{proof}
Orthogonality gives $\|Gx\|_2=\|x\|_2$.  The selected angle annihilates the
$j$th coordinate, so the projection subtracts zero squared magnitude.  The
same argument applies inductively when each subsequent rotation annihilates
the coordinate it is about to discard.
\end{proof}

This lemma is narrower than a general claim about ``noise removal.''  The
operation preserves information from the selected coordinate pair by folding
its norm into another coordinate, but it does not establish that the pair is
noise, preserve direction, preserve attention logits, or improve task reward.
Moreover, the hard projection is algebraically redundant: the adaptive
rotation has already made every projected coordinate zero.  In the audit it
removed only floating-point residue (about $10^{-17}$ of the norm), rather
than the roughly $10^{-3}$ removed by plain truncation.  Norm preservation is
therefore an implementation invariant, not evidence of denoising or useful
information preservation.
In an independent audit over $D\in\{64,128,512,4096\}$ and
$N_{\mathrm{noise}}\in\{1,2,4\}$, maximum relative norm error was
$2.61\times10^{-16}$ in float64 and $1.52\times10^{-8}$ in float32.  The
absolute float32 error stayed below $10^{-7}$ in 11 of 12 configurations and
reached $1.20\times10^{-7}$ at $D=64,N_{\mathrm{noise}}=2$.  The projection
removed roughly $2\times10^{-17}$ of the energy in float64, whereas truncating
the same coordinates would remove $9.6\times10^{-4}$--$6.2\times10^{-2}$ over
the audited grid.  Outside the required $N_{\mathrm{noise}}<D$ contract, the
index map degenerates and measured absolute norm errors exceed $3$; an
implementation must raise rather than continue.

\section{Artifact and Synthetic Stress Tests}

The artifact includes modules for the spectral projection, pairwise score,
Givens elimination, and the modified objective.  Sixteen focused tests pass on
CPU.  They check recurrence values, tensor shapes, nonnegative distance,
rotation invariance, coordinate elimination, fixture construction, output
schema, and a deterministic equal-reward stress test.

The repository harness evaluates $G\in\{4,8,16\}$ and
$L\in\{512,1024,2048\}$ on synthetic log-probability tensors.  The
nondegenerate fixture alternates binary rewards; the homogeneous fixture sets
every reward to one.  Because $A^{\mathrm{aux}}$ is standardized, the magnitude
of the synthetic update is controlled directly by $\lambda$.  Reporting a
single ``recovery percentage'' would therefore report a hyperparameter as if
it were a property of the representation.  Table~\ref{tab:lambda} instead
reports the audit's explicit coupling sweep.

\begin{table}[t]
\centering
\caption{Deterministic float64 CPU measurements as a function of the auxiliary
weight.  The homogeneous column is the gradient created when the reward
gradient is zero.  The cosine is measured against the unmodified gradient on a
heterogeneous control group.  Neither quantity is a learning result.}
\label{tab:lambda}
\small
\begin{tabular}{lccc}
\toprule
$\lambda$ & Homogeneous $\|\nabla\|_2$ & Cosine vs. reward gradient & $1-\cos$ \\
\midrule
$0$       & $0$                  & $1.000000$ & $0$ \\
$10^{-4}$ & $3.735\times10^{-4}$ & $1.000000$ & $1.9\times10^{-9}$ \\
$10^{-2}$ & $3.735\times10^{-2}$ & $0.999981$ & $1.9\times10^{-5}$ \\
$10^{-1}$ & $3.735\times10^{-1}$ & $0.998112$ & $1.9\times10^{-3}$ \\
$1$       & $3.735$              & $0.848370$ & $1.5\times10^{-1}$ \\
\bottomrule
\end{tabular}
\end{table}

The homogeneous gradient is $3.7351\lambda$ across the six audited decades,
exactly the linear behavior Eq.~\eqref{eq:combined-advantage} predicts.  On the
heterogeneous control, increasing auxiliary magnitude progressively rotates
the update away from the reward-driven direction.  A second audit makes the
causal boundary sharper: a forward-path Givens transform has cosine $0.9504$
against the unmodified heterogeneous-group gradient, while the forward-path
spectral transform has cosine $0.2620$.  On homogeneous groups both still
produce exactly zero reward-weighted gradient.  Thus the forward transforms
distort signal where it exists and cannot create it where it does not; only the
explicit auxiliary coupling creates the update.

The repository's separate nine-condition ledger was generated at the default
$\lambda=0.1$.  It also finds identical outputs for the gated and ungated
spectral conditions, so no present artifact demonstrates an incremental Givens
effect.  Its analytic objective-level FLOP model estimates multi-fold rather
than negligible overhead; end-to-end cost remains unmeasured.

\paragraph{What the artifact does not show.}
It contains no language-model fine-tuning with this objective, no held-out
GSM8K or MATH predictions, no multi-seed accuracy comparison, no calibrated
process target, and no evidence that the auxiliary gradient reduces regret.
Synthetic gradient existence is therefore the endpoint of the present
empirical claim.

\section{Prospective Evaluation}

A publishable mitigation claim requires a prospective experiment in which the
trajectory score is defined before outcomes are observed.

\paragraph{Primary hypotheses.}
\textbf{H1 (alignment):} within homogeneous-reward groups, the preregistered
score predicts a later independently measured outcome, such as verifier-
corrected success or pass@k on a disjoint sample.
\textbf{H2 (learning):} adding the score improves fixed held-out accuracy or
sample efficiency at bounded KL divergence and equal rollout/token budget.
\textbf{H3 (gating):} the Givens component adds benefit beyond the same
spectral auxiliary without gating.

\paragraph{Minimum design.}
Use one canonical, completion-masked grpo implementation with clipping and
reference KL; one public model family at $1$--$2$B and $7$--$8$B scales; GSM8K
and a harder MATH subset; and a contamination control.  At both scales and on
both tasks compare: unmodified grpo; spectral auxiliary; Givens-only; combined;
a random score matched in within-group distribution and variance; an
entropy-shaped advantage; a process/verifier-repair signal where available;
and the simple alternatives of dynamic resampling and larger $G$ at equal
tokens.  Give the baseline back the method's extra compute as a budget-matched
control.  Freeze the sampler, reward parser, optimizer, prompts, maximum
length, group size, evaluation set, checkpoint rule, $\lambda$ schedule, and
token budget before inspecting outcomes.

Run the cheapest decisive gate first.  From early, middle, and late frozen
checkpoints at each scale, collect at least $2{,}000$ all-correct and $2{,}000$
all-wrong homogeneous groups.  Without using $s_i$ to construct labels, obtain
verifier-corrected success or disjoint-sample pass@k and test whether $s_i$
predicts it after controlling for completion length and token entropy.  Require
the bootstrap $95\%$ interval for AUC to exclude $0.5$ at both scales and on
both tasks.  If this gate fails, do not train the auxiliary objective; report
the negative alignment result.

If it passes, screen with five paired seeds and expand using
\begin{equation}
 n\ \geq\ \frac{(z_{1-\alpha/2}+z_{1-\beta})^2s_d^2}{\delta^2}+2,
 \label{eq:power}
\end{equation}
where $s_d$ is the standard deviation of within-seed differences and $\delta$
is the preregistered smallest effect worth claiming.  With two-sided
$\alpha=.05$ and $80\%$ power, five seeds can detect only about $1.62s_d$;
smaller effects require expansion.  Correct the four primary contrasts with a
Holm step-down procedure.  Report per-seed learning curves, fixed-budget
held-out accuracy, pass@k, score calibration, reference KL, clip fraction,
output length, GPU-hours, tokens, peak memory, all-correct/all-wrong group
rates, and every failed or collapsed run.

Before that comparison, repair and freeze the implementation contract:
require $1\leq N_{\mathrm{noise}}<D$; build the spectral grid from unmasked
tokens or use a discrete orthogonal basis; normalize entropy as
$H/\log L_{\mathrm{eff}}$; compute the gate per head or token instead of from
a batch-wide mean; and expose $\lambda A_i^{\mathrm{aux}}$ explicitly in the
training loss.  Without the final coupling, a forward-path transform cannot
restore a policy gradient when $A_i^{\mathrm{reward}}=0$.

\paragraph{Falsification criteria.}
Reject the mitigation claim if the score fails out-of-sample alignment, if a
variance-matched random score performs similarly, if gains vanish under paired
seeds, equal tokens, bounded KL, or the budget-matched baseline, if a simple
resampling/larger-group control matches it, or if the Givens ablation adds no
benefit.  A nonzero gradient, norm identity, or passing synthetic test is never
a success criterion.

\section{Limitations and Responsible Claiming}

The auxiliary score is self-referential because it is computed from policy
features or log probabilities.  It may reward novelty, length, formatting, or
instability rather than correctness.  Hidden-state distances are model- and
layer-dependent.  The equispaced Legendre approximation is sensitive to masks
and sequence length.  The adaptive Givens transform is data dependent and may
change directions even when it preserves norms.  Finally, the current stress
tests use idealized tensors and cannot estimate end-to-end memory, latency, or
learning effects.

\section{Conclusion}

Homogeneous rewards remove the centered reward signal.  A trajectory-derived
score can replace that absence with an auxiliary preference, but doing so
changes the objective and creates a new validation burden.  Our prototype has
an executable implementation, an exact continuous identity for its
reconstructed expansion, a norm-preserving adaptive coordinate elimination,
and deterministic synthetic stress tests.  It does not yet have evidence of
better policy learning.  By separating those statements and specifying
falsification criteria, this revision turns an unsupported performance story
into a testable methods proposal.

\appendix
\section{Claim-to-File Index}

Table~\ref{tab:claim-index} maps every numerical empirical statement in the
main text to a machine-readable field.  Audit paths are relative to the
released 77-case audit bundle (SHA-256
\texttt{dbd1386a29d4f413f7a5da56ba6191891915117efa5d13ff6e7eb9ad2c47e5e9});
repository paths are relative to the artifact root.  The independent audit
method passed its internal consistency checks; the earlier manuscript it
audited did not.  These are separate judgments.

\begin{table}[h]
\centering
\caption{Traceability index for the revision's empirical quantities.}
\label{tab:claim-index}
\scriptsize
\begin{tabularx}{\textwidth}{p{.27\textwidth}X}
\toprule
Claim & Artifact pointer \\
\midrule
16 focused repository tests &
\path{pilot/test_spectral_attention.py},
\path{pilot/test_entropic_gating.py}, and
\path{pilot/test_benchmark_spectral_harness.py} \\
Discretization, rank, and $296\times$ trapezoid result &
\path{results/audit_findings_summary.json}:
\path{test2_quadrature.F2.1--F2.5}; raw rows in
\path{results/quadrature_error.csv} and
\path{results/projection_accuracy.csv} \\
$153\%$ padding-grid discrepancy &
\path{results/audit_findings_summary.json}:
\path{test1_invariants} (maximum relative Frobenius discrepancy $1.5268$) \\
Float32/float64 isometry, projection residue, and invalid-width failures &
\path{results/audit_findings_summary.json}:
\path{test3_norms.F3.1--F3.5}; raw sweep in
\path{results/norm_loss_analysis.json} \\
$3.7351\lambda$ sweep and heterogeneous-group cosines &
\path{results/audit_findings_summary.json}:
\path{test4_gradients.F4.3--F4.4}; raw variants in
\path{results/gradient_flow_analysis.json} \\
Binary homogeneous-group floor &
\path{results/audit_findings_summary.json}:
\path{test5_zvf}; raw calculations in
\path{results/zvf_feasibility.json} \\
Default-$\lambda=.1$ nine-condition ledger and FLOP model &
\path{pilot/spectral_benchmark_results.json}:
\path{metadata} and \path{summary}; generator
\path{pilot/benchmark_spectral_harness.py} \\
Prior-art overlaps and citation checks &
\path{results/literature_novelty_matrix.json},
\path{results/verified_literature_matrix.csv}, and
\path{results/citation_live_recheck.json} \\
\bottomrule
\end{tabularx}
\end{table}

\bibliographystyle{plainnat}
\bibliography{flagship}

\end{document}
